\section*{Chapter \ref{ch:rs:cross-sectional-and-cross-asset-features} --- Cross-Sectional and Cross-Asset Features}

\begin{solution}{exo:rs:cross-sectional-and-cross-asset-features:1}
Without itself, the peers' mean is $(1 + 2 + 3 + 4)/4 = 2.5\%$ and the \omterm{def:rs:cross-sectional-and-cross-asset-features:peer}{peer-relative return} is 7.5\%. With itself, the mean is 4\% and
the difference 6\%, which is $\frac45 \times 7.5\%$, as the proposition says with $n = 5$.
\end{solution}

\begin{solution}{exo:rs:cross-sectional-and-cross-asset-features:2}
$10\,000 \times 133/486\,506 = 2.7$; for the top 1\,000, 0.27.
\end{solution}

\begin{solution}{exo:rs:cross-sectional-and-cross-asset-features:3}
$0.1 \times (-6\%) = -0.6\%$ in total, or $-0.6\%/21 = -0.029\%$ a day (2.9 basis points) for 21 days.
\end{solution}

\begin{solution}{exo:rs:cross-sectional-and-cross-asset-features:4}
With $\sigma^2$ the variance rate, $B$'s return over $((k-1)\Delta, k\Delta]$ is $A$'s over $((k-1)\Delta - L, k\Delta - L]$. That interval
overlaps $A$'s previous interval $((k-2)\Delta, (k-1)\Delta]$ over a length $L$ and $A$'s same interval over $\Delta - L$. Brownian increments
over disjoint intervals are independent, so the covariances are $\sigma^2 L$ and $\sigma^2(\Delta - L)$, and dividing by the variance
$\sigma^2\Delta$ of each return gives $L/\Delta$ and $1 - L/\Delta$.
\end{solution}

\begin{solution}{exo:rs:cross-sectional-and-cross-asset-features:5}
Returns on a grid finer than the time between trades are mostly zero, and their correlations shrink (the Epps effect of Book~4,
chapter~21); a grid coarser than the lead blurs it into the same interval. The Hayashi--Yoshida sum uses every observed return and
pairs those whose intervals overlap, so shifting one clock by $\vartheta$ moves the pairs exactly, whatever the spacing of the trades.
\end{solution}

\begin{solution}{exo:rs:cross-sectional-and-cross-asset-features:6}
Regress the small quintile's week on its own previous week and the large quintile's. The two regressors are correlated (the quintiles
move together within a week), so much of what the large quintile's last week seems to predict is the small quintile's own last week. In
the chapter's regression the large quintile's coefficient has a $t$ statistic of only 1.9 from 1962 to 1987, $-0.7$ from 1988 to 2006
and 1.1 since.
\end{solution}

\begin{solution}{exo:rs:cross-sectional-and-cross-asset-features:7}
With a 5-day diffusion the monthly feature's IC falls to 0.008 ($t = 1.4$) while the truth's is 0.082: by the month-end most of the
customer's news has already reached the supplier, and the feature looks at the wrong month. With 63 days, the feature's IC is 0.031
($t = 4.8$) and the truth's 0.050 (the planted drift per month is smaller). A feature's window must match the time the information takes
to diffuse, which the event study of \cref{fig:rs:cross-sectional-and-cross-asset-features:event} measures directly.
\end{solution}

\begin{solution}{exo:rs:cross-sectional-and-cross-asset-features:8}
Tokyo closes before New York opens: the S\&P 500's return of the same date is not known at the Tokyo close, so the correlation is a
leak of timestamps (chapter~3). The tradable version regresses Tokyo's next session on the S\&P's last completed session, known before
Tokyo opens, and must beat Tokyo's own past and the futures that trade overnight.
\end{solution}

\begin{solution}{pb:rs:cross-sectional-and-cross-asset-features:1}
\begin{enumerate}
\item 0.65 times a second in the fast pair, 0.16 in the slow one (about every six seconds).
\item The simulated mids move in steps and trend over a few seconds, so tick-level returns are autocorrelated and their sum of squares
misstates the variance; one-minute returns are long enough to be nearly uncorrelated.
\item At 0.15 seconds: the function is flat within about a second of the latency, because each mid follows its efficient price with a
spread of delays, and the argmax falls anywhere on the flat top.
\item 0.5 seconds. It assumes the two instruments react to their efficient prices with the same distribution of delays, which makes the
cross-correlation symmetric about the latency.
\item 1.17 at 0.5 seconds; 1.03 at zero latency.
\item The centre: 0.053 seconds on average, 0.2 at worst; the argmax: 0.19 on average, 0.5 at worst.
\item Their mids change every six seconds; a shift of a fraction of a second barely changes which intervals overlap, so the function is
flat over seconds and its top is noise.
\item Instruments whose prices change many times within 50 milliseconds, timestamps from one clock at the point of receipt, precise to a
small fraction of that, and enough price changes for the function's top to stand above its noise.
\item 0.59, 0.45 and 0.44.
\item The simulator's mids reach a new efficient price in steps over a few seconds, so every mid-price return trends at short intervals:
each series predicts itself and the other.
\item The first row minus the second: 0.13 at two seconds, 0.17 at five, 0.10 at ten, 0.02 at thirty.
\item At thirty seconds the half-second lead is a small part of each interval and the trends have run their course within it; at 0.1
seconds most intervals contain no price change.
\item In both, the follower's own past explains much of what the leader's past seems to; the test is the leader's contribution beyond it.
\item Probably not with market orders: a one-tick spread is large against a move of a fraction of a tick predicted with correlation 0.6,
and Huth and Abergel found that a naive market-order strategy could not profit from the lead because of the spread. The feature is used to
move quotes or to time passive orders.
\item IC 0.035 with $t = 5.6$; the truth's IC is 0.074.
\item The customer's last month mixes news that has already reached the supplier with news still diffusing, and adds the customer's
market and factor moves; the truth is exactly the part still to come.
\item Daily: 3; monthly: 8; at random: 2.7.
\item \textbf{Named result.} On the fast pair the symmetric centre recovers planted latencies of 0 to 2 seconds within 0.053 seconds on
average (0.2 at worst); the predictability the lead adds over the reverse direction is 0.13 at two seconds, 0.17 at five, 0.10 at ten
and nothing at thirty.
\item The half-second lead: the pair is known, and an hour holds thousands of price changes that all bear on one parameter. The monthly
link hides among 486\,506 candidate pairs with 120 months each; without the customer list it is not found.
\item It predicts the follower beyond the follower's own past, out of sample, at a clock where it survives the cost of trading on it,
with a mechanism that says why.
\end{enumerate}
\end{solution}

\begin{solution}{iq:rs:cross-sectional-and-cross-asset-features:1}
By industry classification, by business description, or by the nearest neighbours in return correlation over a past window (as known at
the time). Pitfalls: including the stock in its own peer mean (a factor $(n-1)/n$ that differs across group sizes), groups too small to
have a stable mean, peers that are not listed on the day (survivorship), and classifications revised after the fact.
\end{solution}

\begin{solution}{iq:rs:cross-sectional-and-cross-asset-features:2}
Take both quote streams on one clock at the point of receipt, compute the Hayashi--Yoshida cross-covariance of the mids with the ETF's
clock shifted over a grid of lags, and take the lag of the peak, or better the centre of symmetry of the function; check it on
simulated data with a known delay, by time of day, and against the time between price changes, which bounds the resolution.
\end{solution}

\begin{solution}{iq:rs:cross-sectional-and-cross-asset-features:3}
Only if the large stocks' past predicts the small stocks beyond their own past. In the French size quintiles, the small quintile's weekly
correlation with its own last week (0.37 from 1962 to 1987) was as large as with the large quintile's (0.31), and controlling for it
leaves the large quintile a $t$ statistic of 1.9; since 2007 the correlation is 0.06.
\end{solution}

\begin{solution}{iq:rs:cross-sectional-and-cross-asset-features:4}
Mostly noise. Among 1\,000 stocks there are 999\,000 ordered pairs; with a few years of data a true link's lagged correlation is a
fraction of the cross-pair standard deviation, so the top 100 are the largest noise draws. On the chapter's synthetic market the top
10\,000 of 486\,506 pairs held 8 of 133 planted links, against 2.7 at random. Start from \omterm{def:rs:cross-sectional-and-cross-asset-features:leadlag}{economic links}, and treat mining as a
multiple-testing problem.
\end{solution}

\begin{solution}{iq:rs:cross-sectional-and-cross-asset-features:5}
Investors with limited attention follow the stock, not its customers (Cohen and Frazzini); information about related firms diffuses
gradually across specialised investors (Menzly and Ozbas). Test with a point-in-time customer list: sort suppliers on their customers'
past-month return, measure the next month's return spread and its $t$ statistic, and check the event-time response of suppliers to
customers' news.
\end{solution}

\begin{solution}{iq:rs:cross-sectional-and-cross-asset-features:6}
If the follower is the leader delayed by $L$ and prices are Brownian, the follower's return over an interval $\Delta \ge L$ has a covariance
$\sigma^2 L$ with the leader's previous interval and $\sigma^2(\Delta - L)$ with the same one: the predictable share is $L/\Delta$, and the
lagged correlation is $L/\Delta$. For $\Delta < L$ all of the covariance lies in past intervals, spread over several. Noise and the
follower's own autocorrelation change the numbers in practice, which is why the chapter measures the asymmetry between the two
directions.
\end{solution}
